\documentclass[11pt]{article}
\usepackage[utf8]{inputenc}
\usepackage{lmodern}
\usepackage[T1]{fontenc}
\usepackage{amsmath,amssymb}
\usepackage[margin=1in]{geometry}
\usepackage{xcolor}
\usepackage{booktabs}
\usepackage{enumitem}
\usepackage[hidelinks]{hyperref}
\setlength{\parindent}{0pt}
\setlength{\parskip}{0.6\baselineskip}

\definecolor{pyblockgreen}{RGB}{0,150,60}
\definecolor{pyblockdark}{RGB}{20,20,20}

\title{\vspace{-1.2cm}\textbf{\color{pyblockgreen}CHIRP}\\[2pt]
\large Coinbase Hashrate-Indexed Reward Payouts\\[4pt]
\normalsize A non-custodial, weighted-lottery shared-reward scheme for Bitcoin mining}
\author{\textit{A PyBLOCK White Paper}}
\date{v1.0 --- June 2026}

\begin{document}
\maketitle

\begin{quote}\itshape
``Without power there is no gravitational wave, and without time there is no spacetime.''\\
\normalfont A gravitational wave needs an energetic source \emph{and} spacetime to travel through.
CHIRP needs the same two things from you: \textbf{power} and \textbf{time}.
\end{quote}

\section*{Abstract}
CHIRP is a Bitcoin mining-pool reward scheme in which a found block's \textbf{coinbase transaction
itself} pays a set of contributing miners directly --- atomically, on-chain, and \textbf{without
the pool ever holding a single satoshi of anyone's earnings}. Because the coinbase output limit
admits only finitely many recipients, CHIRP fills those slots with a \textbf{verifiable weighted
lottery} whose tickets are earned by two things at once: \textbf{tenure} (how long you have mined
with PyBLOCK) and \textbf{power} (the work you have contributed in a rolling window). Over many
blocks, a miner's expected earnings converge to their weight; in any single block, who gets paid is
a publicly auditable draw seeded by the previous block hash.

\section{The dilemma every small pool faces}
A shared mining pool can only ever pay its members \textbf{when it finds a block}. The alternative
--- FPPS, paying per share regardless of luck --- forces the operator to \textbf{front every payout
from reserves} and absorb all variance, which is bankruptcy waiting to happen for anyone without
warehouse-scale hashrate.

The honest scheme at small scale is therefore some form of pay-per-last-$N$-shares: rare, large,
fairly-split payouts. But PPLNS as usually deployed has a second problem --- \textbf{custody}. The
block reward lands in the operator's wallet, and the operator then distributes it. CHIRP removes
the custody problem entirely.

\section{Prior art}
\begin{itemize}[leftmargin=1.4em]
  \item \textbf{PPS / FPPS} --- predictable pay, operator-fronted, variance on the operator. Needs
  scale and reserves.
  \item \textbf{PPLNS} --- pay over a window of recent shares. Fair, but typically custodial.
  \item \textbf{OCEAN's TIDES} --- a PPLNS variant over a rolling window of the last 8 blocks
  ($8 \times D_{\text{net}}$ of share-work) that pays each miner directly in the coinbase ---
  non-custodial and auditable. TIDES is the closest prior art and the inspiration for CHIRP.
\end{itemize}

\section{CHIRP in one sentence}
\begin{quote}
When the pool solves a block, its coinbase pays a weighted-lottery sample of eligible miners
directly on-chain, where each miner's odds and slice are the average of how long they have mined
here (\textbf{days}) and how much work they have recently contributed (\textbf{power}).
\end{quote}

\section{The mechanism}

\subsection{Eligibility --- the candidacy gate}
A miner is an \emph{aspirant} only if they clear \emph{both} floors:
\[
  \text{candidate} \iff \big( \text{days}_i \ge \text{MIN\_DAYS} \big) \;\wedge\;
  \big( \text{power}_i \ge \text{MIN\_POWER} \big).
\]
\textsc{min\_power} eliminates dust; \textsc{min\_days} sets a loyalty floor. \emph{Whoever fails
either floor is not an aspirant}: a CPU idling ``connected for years'' never clears \textsc{min\_power}, and a whale
that parachutes in never clears \textsc{min\_days}.

\subsection{Weight --- the average of both}
Among candidates, each miner's weight is the \textbf{mean} of two normalized factors:
\[
  w_i \;=\; \frac{\,d_i + p_i\,}{2},
  \qquad
  d_i = \min\!\Big(\frac{\text{days}_i}{\text{DAYS\_FULL}},\, 1\Big),
  \qquad
  p_i = \min\!\Big(\frac{\text{power}_i}{\text{POWER\_FULL}},\, 1\Big).
\]
The mean --- not a product, not a sum --- is deliberate: it \textbf{balances} the inputs so neither
alone dominates. A veteran with a weak machine is dragged down by low $p_i$; a freshly-joined giant
is dragged down by low $d_i$. To win you need \textbf{both}. Here $\text{power}_i$ is the miner's
accumulated share-work (the sum of share difficulties) inside the rolling window.

\subsection{The lottery --- filling the coinbase slots}
A coinbase can carry only $\text{MAX\_N}$ outputs before competing with fee-paying transactions for
block weight. When candidates exceed slots, CHIRP \textbf{draws} winners with probability
proportional to weight, using the Efraimidis--Spirakis weighted-reservoir method:
\[
  \text{seed} = \text{prevBlockHash}_{64},
  \qquad
  \text{key}_i = u_i^{\,1/w_i},
  \qquad
  u_i = U(\text{seed}, \text{addr}_i) \in (0,1),
\]
\[
  \text{winners} = \operatorname*{top\text{-}MAX\_N}_{i}\; \text{key}_i .
\]
The uniforms $u_i$ are derived deterministically from the seed and the miner's address. The seed is
the \textbf{previous block hash}: unknown when miners began contributing (so it cannot be gamed),
yet fully reproducible afterward (so anyone can audit the draw). \textbf{The coinbase-size constraint
thus becomes the mechanism} --- a fair, recurring lottery for the slots.

\subsection{Payout --- the split among winners}
The reward, less the pool fee, is divided across winners in proportion to weight:
\[
  \text{payout}_i \;=\; \big(R - \text{fee}\big)\cdot
  \frac{w_i}{\displaystyle\sum_{j \in \mathcal{W}} w_j},
  \qquad
  R = \text{subsidy} + \text{tx fees}.
\]
Every payout is a coinbase \texttt{TxOut}. If any miner solves the block, the \emph{network} pays
everyone at once --- no balances, no withdrawals, no custody.

\subsection{The window --- a bounded variant of TIDES}
\textsc{power} is measured over a rolling window of recent share-work. TIDES uses an
$8 \times D_{\text{net}}$ window; at a small pool's hashrate that would span millions of years and
never slide. CHIRP therefore bounds it:
\[
  \text{window} \;=\; \min\!\big(\,8 \times D_{\text{net}},\; 24\text{h}\,\big).
\]
This converges to TIDES at exahash scale while staying meaningful and fair for active miners today.
Because share-work is $\text{hashrate} \times \text{time}$, the window already rewards both
contribution and recent presence; the explicit days/power average layers tenure on top.

\section{The gravitational-wave principle}
CHIRP is named for the \textbf{chirp} --- the rising signal LIGO detects when two black holes spiral
together and merge, radiating ripples through spacetime. The metaphor is exact, and it is also the
design's first principle:
\begin{itemize}[leftmargin=1.4em]
  \item \textbf{Without power, there is no wave.} Gravitational waves are radiated by accelerating
  mass-energy. No source $\Rightarrow$ no signal. No real hashrate $\Rightarrow$ no CHIRP weight.
  \item \textbf{Without time, there is no spacetime.} The wave is a ripple of \emph{space-time};
  remove the temporal dimension and there is no medium to ripple. No tenure $\Rightarrow$ no weight.
\end{itemize}
The dual requirement is not arbitrary --- it is \textbf{cosmological}. Miners \emph{merge} their
hashrate like colliding black holes, and when the block is found the reward \textbf{chirps} out as a
wave that pays every contributor. (Ocean $\to$ TIDES, tides of water. Gravity $\to$ CHIRP, ripples
of spacetime.)

\section{Verifiability \& transparency}
CHIRP is auditable end to end: \textbf{weights} are computed from public inputs (each address's
first-seen timestamp and rolling share-work) and published; \textbf{the draw} is deterministic given
the previous block hash, so anyone can replay it and confirm their odds; \textbf{the split} is
visible directly in the coinbase of every block found. Don't trust --- verify.

\section{Anti-gaming}
\begin{center}
\begin{tabular}{@{}ll@{}}
\toprule
\textbf{Attack} & \textbf{Why it fails} \\
\midrule
Sybil (split across new addresses) & each new address starts at 0 days $\rightarrow$ below \textsc{min\_days} \\
Idle loyalty (a CPU for years) & never clears \textsc{min\_power}; the average sinks low power \\
Whale parachute (power, no tenure) & never clears \textsc{min\_days}; $d_i = 0$ sinks the average \\
Window-boundary hopping & muted by 24h smoothing and the lottery randomness \\
Predicting the draw & the seed is the \emph{previous} block hash --- unknown until it exists \\
\bottomrule
\end{tabular}
\end{center}

\section{Parameters (v1 defaults, tunable)}
\begin{center}
\begin{tabular}{@{}lll@{}}
\toprule
\textbf{Parameter} & \textbf{Default} & \textbf{Role} \\
\midrule
\textsc{min\_days}   & 7 days & loyalty floor for candidacy \\
\textsc{min\_power}  & $\approx 10$ shares at min-difficulty & dust floor for candidacy \\
\textsc{days\_full}  & 30 days & tenure at which $d_i$ saturates \\
\textsc{power\_full} & calibrated to a healthy ASIC / 24h & work at which $p_i$ saturates \\
\textsc{max\_n}      & 100 & coinbase slots = lottery winners \\
Window               & $\min(8 D_{\text{net}}, 24\text{h})$ & rolling power window \\
Fee                  & 0.9\% & PyBLOCK's coinbase cut \\
\bottomrule
\end{tabular}
\end{center}

\section{Deployment}
CHIRP runs as a dedicated Stratum V2 pool instance on port \textbf{5554}, alongside (and entirely
independent of) PyBLOCK's solo/lottery stratum on 5555. Both serve PyBLOCK's own \textbf{BIP-110}
templates from the same Bitcoin Knots node, so every CHIRP block also signals BIP-110. A per-address
stats registry, updated on every accepted share and snapshotted to disk, supplies the tenure and
power that become each miner's weight.

\section{Limitations \& honesty}
\textbf{Variance is real.} A small pool finds blocks rarely; CHIRP smooths \emph{who shares} a
block, not \emph{how often} blocks arrive. It is a syndicate, not a salary --- miners who want
steady income should mine DATUM $\to$ OCEAN, which PyBLOCK also offers. \textsc{power\_full} and
\textsc{min\_power} require calibration against live hashrate.

\section{Conclusion}
CHIRP turns three constraints into features. The \emph{custody problem} becomes \textbf{non-custody}
--- the network pays miners directly in the coinbase. The \emph{coinbase-size limit} becomes a
\textbf{fair, verifiable lottery} for the slots. And the \emph{small-pool variance} becomes the
\textbf{honest soul of the product} --- a syndicate that, by the same logic as the cosmos, pays
those who bring both energy and time.

\begin{center}\itshape
Mine on PyBLOCK. Merge your hashrate. Hear the chirp. \quad $\bigstar$
\end{center}

\vfill
\begin{center}\small
\textit{PyBLOCK --- The CypherPunk's Bitcoin Mining Pool. Powered by Bitcoin Knots $\cdot$ BIP-110 enforced.}
\end{center}

\end{document}
